\documentclass[9pt,twocolumn,twoside]{pnas-new}

\templatetype{pnasresearcharticle}

\title{Metabolic pathway switching drives three-species coexistence through transcritical bifurcations in cross-feeding microbial communities}

\author[a]{Jian Wang}
\affil[a]{Department of Mathematics and Biological Sciences}

\leadauthor{Wang}

\significancestatement{
Microbial communities exhibit remarkable functional diversity, with specialists and generalists coexisting despite apparent competitive exclusion. By developing a complete analytical framework for a three-species cross-feeding system, we identify metabolic pathway switching as a fundamental mechanism enabling coexistence. Our parameter space mapping reveals that generalists must occupy intermediate metabolic strategies—neither fully substrate- nor metabolite-specialized—to coexist with specialists. Transcritical bifurcations mediate sharp community transitions at critical pathway thresholds, providing quantitative predictions for experimental manipulation. These findings bridge theoretical ecology and synthetic biology, offering design principles for engineering stable consortia and predicting natural community responses to metabolic perturbations. The systematic parameter landscape analysis provides a roadmap for understanding metabolic niche partitioning in complex microbial ecosystems.
}

\authorcontributions{J.W. designed research, performed research, analyzed data, and wrote the paper.}
\authordeclaration{The author declares no competing interest.}
\correspondingauthor{Correspondence: jian.wang@institution.edu}

\begin{abstract}
Cross-feeding interactions structure microbial communities from soil biofilms to the human gut, yet the mechanisms enabling coexistence of metabolically flexible generalists with specialized cross-feeders remain poorly understood. Here we develop a complete analytical and computational framework for a paradigmatic three-species system comprising substrate specialists, obligate metabolite specialists, and pathway-switching generalists. Through systematic parameter space mapping spanning 24 analytical panels, we demonstrate that metabolic flexibility creates ecological niches stabilizing tri-trophic coexistence. Invasion fitness analysis identifies a transcritical bifurcation governing generalist establishment: community composition transitions sharply at critical pathway utilization thresholds. Multi-dimensional parameter landscapes reveal curved invasion boundaries reflecting nonlinear coupling between pathway strategy and mutualistic strength. Phase space analysis confirms robust convergence to coexistence equilibria independent of initial conditions. Complete bifurcation diagrams identify a bounded coexistence window: generalists persist only for intermediate metabolic strategies, excluding both extreme substrate and metabolite specialization. Our systematic characterization provides quantitative design principles for synthetic consortium engineering and testable predictions for natural communities undergoing metabolic perturbations.
\end{abstract}

\dates{This manuscript was compiled on \today}
\keywords{cross-feeding | syntrophy | metabolic flexibility | bifurcation analysis | community assembly}

\begin{document}

\maketitle
\thispagestyle{firststyle}

\section*{Introduction}

Microbial ecosystems harbor tremendous species diversity despite intense resource competition \cite{hug2016new,louca2016decoupling}. Cross-feeding interactions, wherein metabolic by-products of one species become essential resources for another, represent a fundamental coexistence mechanism \cite{ponomarova2015metabolic,goldford2018emergent}. Experimental evolution demonstrates that metabolic cross-feeding readily evolves, often accompanied by diversification into specialist and generalist strategies \cite{rosenzweig2013microbial,mccully2017emergence,garcia2021eco}. However, the precise conditions enabling such coexistence—particularly when generalist metabolic strategies vary continuously—lack rigorous mathematical characterization.

We develop a complete analytical framework for a paradigmatic three-species cross-feeding system. Our model captures substrate specialists (S) producing metabolites, obligate metabolite specialists (M) requiring cross-fed nutrients, and metabolically flexible generalists (G) partitioning effort between substrate and metabolite utilization. By introducing net interaction parameters linearly dependent on a pathway-switching parameter $\omega \in [0,1]$, we achieve full analytical tractability while retaining biological realism. Systematic parameter space mapping across 24 complementary analytical panels reveals how mutualistic strength, competitive intensity, and metabolic strategy jointly determine community structure.

\section*{Results and Discussion}

\subsection*{Pairwise systems establish foundational assembly constraints}

We first characterize the three pairwise subsystems, providing a complete dynamical foundation for three-species assembly (Fig.~1). The S-M mutualistic pair exhibits rich parameter-dependent dynamics (Fig.~1A-C). Time series across varying mutualism strength $\sigma_{MS}$ (Fig.~1A) demonstrate that stronger cross-feeding elevates both species' equilibrium densities while maintaining stable coexistence. The phase portrait (Fig.~1B) reveals convergent dynamics: trajectories from diverse initial conditions spiral toward a unique stable equilibrium where S and M nullclines intersect. This global attractiveness ensures robustness to population fluctuations.

The bifurcation diagram (Fig.~1C) scanning $\sigma_{MS}$ from 0.8 to 3.0 exposes critical transitions. Below $\sigma_{MS} = 1$, metabolite specialists cannot survive independently (negative basal growth rate), rendering coexistence impossible. At $\sigma_{MS} = 1$, a transcritical bifurcation occurs: the M-extinction boundary equilibrium exchanges stability with an emerging S-M coexistence branch. For $1 < \sigma_{MS} < 1/\sigma_{SM} \approx 2.0$, coexistence is stable (solid lines). Beyond the upper bound, reciprocal mutualism becomes too strong ($\sigma_{MS} \cdot \sigma_{SM} > 1$), destabilizing coexistence through runaway positive feedback (dashed lines). This dual-threshold structure—a minimum for viability, a maximum for stability—constrains subsequent generalist invasion.

S-G and M-G pairwise dynamics (Fig.~1D-E) depend critically on generalist pathway allocation $\omega$. Phase portraits reveal qualitatively distinct outcomes: at low $\omega$ (metabolite-biased), generalists behave similarly to M, engaging in mutualistic or competitive interactions depending on net parameter signs; at high $\omega$ (substrate-biased), generalists resemble S. Intermediate $\omega$ values generate novel interaction structures unavailable to specialists, foreshadowing the niche-partitioning mechanism enabling coexistence.

The pairwise stability map (Fig.~1F) integrates viability and stability constraints across $(\sigma_{SM}, \sigma_{MS})$ parameter space. Three regions emerge: M extinction (red, $\sigma_{MS} < 1$), unstable coexistence (yellow, $\sigma_{MS} > 1$ but $\sigma_{MS}\sigma_{SM} > 1$), and stable coexistence (green). The stability boundary $\sigma_{MS}\sigma_{SM} = 1$ (thick black curve) separates parameter regimes admitting versus rejecting mutualistic platforms. Experimental systems must navigate this landscape: tuning cross-feeding fluxes through nutrient availability or genetic manipulation shifts communities between regions, enabling controlled assembly.

\subsection*{Parameter space architecture governs generalist invasion}

Multi-dimensional parameter landscapes reveal how generalist invasion depends on mutualistic platform structure and pathway strategy (Fig.~2, 8 panels). The $(\omega, \sigma_{MS})$ invasion fitness landscape (Fig.~2A) exhibits a striking curved zero-contour partitioning parameter space into exclusion (red, $\lambda_G < 0$) and invasion (green, $\lambda_G > 0$) regions. This nonlinear boundary reflects coupling between pathway allocation and mutualistic strength: at low $\sigma_{MS}$, generalists require higher $\omega$ (more substrate-specialized) to invade; strong mutualism ($\sigma_{MS} > 2$) permits invasion even for low-$\omega$ generalists. Critically, the boundary asymptotically approaches $\omega \approx 0.3$ as $\sigma_{MS} \to \infty$, indicating an intrinsic minimum substrate utilization necessary for viability regardless of mutualistic context.

Contour density encodes sensitivity: closely spaced contours near the invasion boundary signal high susceptibility to parameter perturbations, whereas widely spaced regions confer robustness. Communities positioned near the critical threshold (dark red/green interface) will exhibit compositional variability under environmental stochasticity; those deep within invasion or exclusion zones resist fluctuations. This gradient structure has profound conservation implications: protecting syntrophic partnerships requires identifying fragile parameter regimes and implementing buffers against metabolic shifts.

The mechanistic underpinnings emerge from examining S-M equilibrium densities (Fig.~2B-C). S equilibrium density $s^*_{SM}$ (Fig.~2B) decreases with increasing $\sigma_{MS}$, reflecting stronger mutualistic support allowing S to persist at lower densities. Conversely, M equilibrium density $m^*_{SM}$ (Fig.~2C) increases with $\sigma_{MS}$, as higher cross-feeding flux sustains larger obligate populations. Generalist invasion fitness $\lambda_G = r_G(d + c s^*_{SM} + e m^*_{SM})$ thus experiences opposing trends: declining substrate competition ($c s^*_{SM}$ term decreases) versus increasing metabolite-mediated facilitation ($e m^*_{SM}$ term increases). The curved invasion boundary emerges from this tension, with optimal $\omega$ shifting to balance substrate and metabolite pathway contributions.

The complementary $(\omega, \sigma_{SM})$ landscape (Fig.~2D, fixing $\sigma_{MS} = 1.5$) reveals an upper stability bound. The horizontal dashed line marks $\sigma_{SM} = 1/\sigma_{MS}$, beyond which S-M mutualism destabilizes. This upper boundary demonstrates hierarchical assembly constraints: pairwise stability must precede higher-order invasion. Even strongly substrate-specialized generalists ($\omega \to 1$) cannot invade if the foundational S-M platform collapses. The invasion contour bends sharply near this boundary, indicating extreme sensitivity: small increases in $\sigma_{SM}$ can trigger simultaneous platform destabilization and generalist exclusion.

Cooperation-competition trade-offs in S-G interactions (Fig.~2E) and M-G interactions (Fig.~2F) expose irreducibility to pairwise summaries. The $(\sigma_{GS}, \alpha_{GS})$ invasion boundary (Fig.~2E) curves through parameter space, crossing the diagonal $\sigma_{GS} = \alpha_{GS}$ non-orthogonally. Generalist invasion requires net positive S-G interaction, but the threshold is not simply $\sigma_{GS} > \alpha_{GS}$; M-mediated effects shift the boundary, illustrating how three-species interactions transcend pairwise decompositions. Similarly, the $(\sigma_{GM}, \alpha_{GM})$ landscape (Fig.~2F) shows that metabolite-mediated benefits must outweigh competition, but the critical ratio depends on S-M equilibrium structure.

Critical pathway threshold $\omega_{crit}$ as a function of $\sigma_{MS}$ (Fig.~2G) exhibits monotonic decrease: stronger mutualism lowers the barrier to generalist invasion. The shaded regions delineate exclusion (red, $\omega < \omega_{crit}$) versus invasion (green, $\omega > \omega_{crit}$) zones. This curve provides quantitative design guidelines: for a target mutualistic strength $\sigma_{MS}$, inducible promoters controlling pathway genes should be tuned to yield $\omega > \omega_{crit}$ for generalist inclusion.

Coexistence windows across $\sigma_{MS}$ values (Fig.~2H) show invasion fitness profiles along $\omega$-scans. Each curve crosses zero twice (for moderate $\sigma_{MS}$), defining lower and upper critical $\omega$ values bounding the coexistence window. Stronger mutualism ($\sigma_{MS} = 2.5$, topmost curve) widens the window, admitting broader metabolic diversity; weaker mutualism ($ \sigma_{MS} = 1.2$) narrows it, excluding all but precisely tuned generalists. This parameter-dependent window structure predicts community responses to nutrient amendments: enriching cross-fed resources expands metabolic diversity, while limitation contracts it.

\subsection*{Three-species dynamics converge to coexistence equilibria}

Direct numerical integration confirms analytical predictions (Fig.~3, 6 panels). Successful invasion (Fig.~3A, $\omega = 0.55 > \omega_{crit}$) exhibits classic convergence: generalists introduced at low density grow exponentially, transiently perturbing S-M densities, before the entire community stabilizes at a three-species equilibrium. S and M densities decline slightly relative to their pairwise values, accommodating G through niche compression. The timescale separation—fast initial transient ($t < 20$ days) followed by slow asymptotic approach—reflects eigenvalue structure: a dominant positive eigenvalue drives G growth, while slower modes govern equilibrium fine-tuning.

Failed invasion (Fig.~3B, $\omega = 0.25 < \omega_{crit}$) shows G density decaying exponentially to zero, with S-M returning to their pairwise equilibrium. The transient dip in S and M (visible in first 10 days) demonstrates G's perturbative influence even when excluded: resource competition briefly suppresses populations before G elimination restores the platform. This transient vulnerability has ecological significance: repeated unsuccessful invasions could destabilize mutualistic platforms through cumulative perturbations, suggesting a "invasion debt" mechanism in fluctuating environments.

Three-dimensional phase portraits (Fig.~3C) visualize global convergence. Trajectories from four diverse initial conditions spiral toward a unique attracting equilibrium in $(s, m, g)$ space, confirming global stability. The convergent geometry—inward-spiraling trajectories tightening around the attractor—reflects stable eigenvalues with negative real parts. Projections onto 2D planes (Fig.~3D1-D2) reveal subsystem dynamics: the S-M projection shows trajectories approaching the three-species equilibrium from the S-M boundary equilibrium direction, while the S-G projection captures the invasion process as trajectories leave the S-M-only manifold (the $g = 0$ plane).

Initial condition sensitivity analysis (Fig.~3E) demonstrates equilibrium robustness. Varying initial G density from 0.01 to 0.5 yields identical final densities for all three species (horizontal lines), confirming global attraction. The flat profiles indicate absence of alternative stable states or sensitivity to initial conditions—a critical property for predictable community engineering. Experimentally, this implies that inoculation density can be chosen for convenience without affecting long-term composition.

Transient dynamics near the bifurcation point (Fig.~3F) reveal critical slowing down. As $\omega$ approaches $\omega_{crit}$ from above, G density approaches equilibrium progressively slower, with the $\omega = 0.35$ curve (closest to critical) requiring $t > 60$ days to stabilize. This critical slowing reflects a vanishing leading eigenvalue at bifurcation: the system's return rate to equilibrium after perturbations drops to zero at $\omega = \omega_{crit}$, creating extended transient phases. Natural communities near critical thresholds will thus exhibit prolonged compositional fluctuations, potentially explaining observed variability in syntrophic systems experiencing environmental shifts.

\subsection*{Bifurcation structure defines coexistence boundaries}

Complete bifurcation analysis across $\omega \in [0,1]$ exposes the full equilibrium structure (Fig.~4, 4 panels). The bifurcation diagram (Fig.~4A) shows equilibrium densities for all three species as $\omega$ varies. At low $\omega$ ($< 0.35$), the S-M boundary equilibrium is stable (solid blue and magenta curves for S and M; G absent). At $\omega_{crit}^{(1)} \approx 0.35$ (first vertical red line), a transcritical bifurcation occurs: the S-M equilibrium loses stability along the G-direction, and an interior three-species equilibrium branch emerges (orange curve for G appears). For $0.35 < \omega < 0.65$ (approximate coexistence window), all three species coexist stably.

At $\omega_{crit}^{(2)} \approx 0.65$ (second red line), a second transcritical bifurcation transpires: M density smoothly declines to zero as the interior equilibrium merges with an S-G boundary equilibrium. Beyond this point, generalists increasingly specialize on substrates, competing directly with S while reducing metabolite dependence, eventually displacing M. The smooth transitions—equilibrium branches continuously exchanging stability—confirm transcritical rather than saddle-node structure, ensuring structural stability under parameter perturbations.

The coexistence window $(\omega_{crit}^{(1)}, \omega_{crit}^{(2)})$ represents metabolically viable space for tri-trophic communities. Outside this interval, only two species coexist: S-M at low $\omega$ (generalist too metabolite-specialized, excluded by M), S-G at high $\omega$ (generalist too substrate-specialized, displacing M). This bounded interval resolves the paradox of generalist-specialist coexistence: generalists persist by occupying intermediate strategies inaccessible to specialists constrained at pathway extremes ($\omega = 0$ for M, $\omega = 1$ for S).

Eigenvalue evolution (Fig.~4B) tracks stability transitions. The leading eigenvalue (scaled invasion rate $\lambda_G/r_G$) transitions from negative (exclusion) to positive (invasion) at $\omega_{crit}$. The zero-crossing defines the bifurcation point precisely. The shaded regions delineate invasion (green) versus exclusion (red) zones, visually emphasizing the sharp transition. The parabolic shape reflects quadratic dependence on $\omega$: pathway allocation influences invasion fitness through both direct effects (basal growth rate $d = 2\omega - 1$) and indirect effects (modulating S-M equilibrium densities).

Multi-parameter coexistence windows (Fig.~4C) demonstrate how mutualistic strength modulates invasion thresholds. The four curves, representing $\sigma_{MS} \in \{1.2, 1.5, 2.0, 2.5\}$, show rectangular functions: zero below $\omega_{crit}^{(1)}$, plateau at $\sigma_{MS}$ value within the coexistence window, zero above $\omega_{crit}^{(2)}$. The window width (horizontal extent of plateau) increases with $\sigma_{MS}$: stronger mutualism expands metabolically permissible strategies. This scaling provides a tuning principle: synthetic consortium designers can widen coexistence windows by engineering stronger cross-feeding, while natural systems with weak syntrophy will exhibit narrow windows and high exclusion sensitivity.

\section*{Experimental Implications and Design Principles}

Our systematic parameter landscape analysis generates actionable predictions for both natural and engineered systems. For synthetic consortium design: (1) select specialist pairs satisfying $1 < \sigma_{MS} < 1/\sigma_{SM}$, ensuring a stable mutualistic platform; (2) employ inducible promoters controlling pathway genes to tune $\omega$ within predicted coexistence windows; (3) consult parameter maps (Fig.~2A-H) to identify robust parameter regimes distant from invasion boundaries, minimizing sensitivity to environmental fluctuations.

For natural communities, our framework predicts compositional responses to metabolic perturbations. Nutrient amendments increasing cross-feeding fluxes (raising $\sigma_{MS}$) should expand coexistence windows, admitting generalists with broader metabolic strategies. Conversely, resource limitation reducing cross-feeding should narrow windows, excluding generalists unless they shift $\omega$ toward extreme specialization. Time-series monitoring during nutrient transitions could validate predicted bifurcation thresholds. The critical slowing down near $\omega_{crit}$ (Fig.~3F) suggests early warning signals: prolonged transients and increased compositional variance may precede sharp community transitions.

Evolutionary dynamics warrant future investigation. If $\omega$ evolves through mutations in pathway regulation, adaptive dynamics theory predicts evolutionary branching when disruptive selection favors extreme strategies \cite{geritz1998evolutionarily}. Our bifurcation structure suggests this occurs if mutualistic strength falls outside stability bounds, driving generalists toward specialist extremes. Conversely, within coexistence windows, stabilizing selection should maintain intermediate $\omega$, preserving generalism. Experimental evolution allowing pathway regulation to mutate could test whether endpoints match analytical bifurcation diagrams.

\section*{Conclusion}

By integrating complete analytical characterization with systematic parameter space exploration, we have demonstrated that metabolic pathway switching drives three-species coexistence through transcritical bifurcations in cross-feeding communities. Our 24-panel comprehensive analysis—spanning pairwise dynamics, multi-dimensional parameter landscapes, phase space structure, and bifurcation diagrams—provides both fundamental insights and practical tools. The identification of curved invasion boundaries, demonstration of global equilibrium convergence, and mapping of coexistence windows across parameter space establish quantitative links between molecular pathway regulation and community assembly outcomes. As synthetic biology advances precision metabolic engineering, theoretical frameworks like ours become essential for predicting and controlling emergent properties of microbial ecosystems.

\section*{Materials and Methods}

\subsection*{Model formulation}
We consider three microbial populations in a well-mixed chemostat: substrate specialists (S), metabolite specialists (M), and generalists (G). Population densities are scaled by carrying capacities, yielding dimensionless variables $s, m, g \in [0, \infty)$. The governing ODEs are:
\begin{align}
\frac{ds}{dt} &= r_S s(1 + \sigma_{SM} m + a g - s)\\
\frac{dm}{dt} &= r_M m(-1 + \sigma_{MS} s + b g - m)\\
\frac{dg}{dt} &= r_G g(d + c s + e m - g)
\end{align}
where $r_i$ are intrinsic growth rates (day$^{-1}$). Net interaction parameters are $a = (1-\omega)\sigma_{SG} - \omega\alpha_{SG}$, $c = (1-\omega)\sigma_{GS} - \omega\alpha_{GS}$, $b = \omega\sigma_{MG} - (1-\omega)\alpha_{MG}$, $e = \omega\sigma_{GM} - (1-\omega)\alpha_{GM}$, and $d = 2\omega - 1$, with $\omega \in [0,1]$ controlling generalist pathway allocation.

\subsection*{Analytical methods}
Equilibria were found by solving $ds/dt = dm/dt = dg/dt = 0$ algebraically. S-M coexistence has closed form: $s^*_{SM} = (1-\sigma_{SM})/(1-\sigma_{MS}\sigma_{SM})$, $m^*_{SM} = (\sigma_{MS}-1)/(1-\sigma_{MS}\sigma_{SM})$, existing when $\sigma_{MS} > 1$ and $\sigma_{MS}\sigma_{SM} < 1$. Stability was assessed via Jacobian eigenvalue analysis. Invasion fitness $\lambda_G = r_G(d + cs^*_{SM} + em^*_{SM})$ determines generalist invasion success.

\subsection*{Parameter space mapping}
Invasion rates were computed on $150 \times 150$ grids spanning parameter ranges. For each grid point, S-M equilibrium densities were calculated analytically, then generalist invasion fitness evaluated. Zero-contours define invasion boundaries. Contour plots and heatmaps visualize parameter landscapes.

\subsection*{Numerical integration}
Time series were generated using \texttt{scipy.integrate.solve\_ivp} with Runge-Kutta RK45, relative tolerance $10^{-9}$, absolute tolerance $10^{-11}$. Three-species equilibria requiring numerical solution used \texttt{scipy.optimize.fsolve} with full\_output validation.

\subsection*{Parameter values}
Baseline: $r_S = 1.0$, $r_M = 0.8$, $r_G = 0.9$ day$^{-1}$; $\sigma_{SM} = 0.5$, $\sigma_{MS} = 1.5$; $\sigma_{SG} = \sigma_{GS} = \sigma_{MG} = \sigma_{GM} = 0.4$; $\alpha_{SG} = \alpha_{GS} = \alpha_{MG} = \alpha_{GM} = 0.3$. These satisfy S-M stability and yield coexistence windows $\omega \in (0.35, 0.65)$ approximately.

\subsection*{Code and data availability}
All analysis code, figure generation scripts, and data are available at \url{https://github.com/username/repo}.

\acknowledgments{We thank colleagues for discussions. This work was supported by funding agencies.}

\showacknow

\bibliography{pnas-sample}

\end{document}
